\chapter{Order continuity}

This chapter studies normed vector lattices whose norm is compatible
with order convergence. A norm is \emph{$\sigma$-order continuous}
when every antitone sequence of positive elements with infimum $0$
converges to $0$ in norm, and \emph{order continuous} when the
analogous condition holds for nets. After recording the basic
sequential consequences of $\sigma$-order continuity and the fact
that order continuity forces Dedekind completeness in the Banach
lattice setting, we prove three landmark equivalent
characterisations of order continuity for a Banach lattice:
Nakano's theorem (in terms of monotone bounded sequences),
Meyer--Nieberg's theorem (in terms of order bounded disjoint
sequences), and Ando's theorem (in terms of norm-closed ideals).
We close with the decomposition lemma: in an order continuous
Banach lattice, the principal band projections along a maximal
disjoint family form an unconditional decomposition of every
element.

\section{Order continuous norms}

This section introduces the two classes of order continuous norms,
records that the order continuous version is stronger, and
establishes the equivalent sequential characterisations of
$\sigma$-order continuity. We then show that an order continuous
Banach lattice is automatically Dedekind complete.

\begin{definition}
\label{def:isSigmaOrderContinuousNorm}
\lean{IsSigmaOrderContinuousNorm}
\leanok
\uses{def:normed-vector-lattice}
Let $X$ be a normed vector lattice. The norm of $X$ is
\emph{$\sigma$-order continuous} if, for every antitone sequence
$(u_n)_{n \in \mathbb{N}}$ in $X$ with $0 \le u_n$ for all $n$ and
greatest lower bound $0$, the sequence $u_n$ converges to $0$ in
norm.
\end{definition}

\begin{definition}
\label{def:isOrderContinuousNorm}
\lean{IsOrderContinuousNorm}
\leanok
\uses{def:normed-vector-lattice}
Let $X$ be a normed vector lattice. The norm of $X$ is
\emph{order continuous} if, for every nonempty upward-directed
preorder $\iota$ and every antitone net $(u_i)_{i \in \iota}$ in
$X$ with $0 \le u_i$ for all $i$ and greatest lower bound $0$, the
net $u_i$ converges to $0$ in norm.
\end{definition}

\begin{lemma}
\label{lem:isOrderContinuousNorm-toIsSigmaOrderContinuousNorm}
\lean{IsOrderContinuousNorm.toIsSigmaOrderContinuousNorm}
\leanok
\uses{def:isSigmaOrderContinuousNorm, def:isOrderContinuousNorm}
Let $X$ be a normed vector lattice with order continuous norm.
Then the norm of $X$ is $\sigma$-order continuous.
\end{lemma}

\begin{lemma}
\label{lem:isSigmaOrderContinuousNorm-tendsto-of-monotone-isLUB}
\lean{IsSigmaOrderContinuousNorm.tendsto_of_monotone_isLUB}
\leanok
\uses{def:isSigmaOrderContinuousNorm}
Let $X$ be a normed vector lattice with $\sigma$-order continuous
norm. If $(u_n)_{n \in \mathbb{N}}$ is a monotone sequence in $X$
with least upper bound $x \in X$, then $u_n$ converges to $x$ in
norm.
\end{lemma}

\begin{lemma}
\label{lem:isSigmaOrderContinuousNorm-tendsto-of-antitone-isGLB}
\lean{IsSigmaOrderContinuousNorm.tendsto_of_antitone_isGLB}
\leanok
\uses{def:isSigmaOrderContinuousNorm}
Let $X$ be a normed vector lattice with $\sigma$-order continuous
norm. If $(u_n)_{n \in \mathbb{N}}$ is an antitone sequence in $X$
with greatest lower bound $x \in X$, then $u_n$ converges to $x$
in norm.
\end{lemma}

\begin{lemma}
\label{lem:isSigmaOrderContinuousNorm-tendsto-of-abs-sub-le-antitone}
\lean{IsSigmaOrderContinuousNorm.tendsto_of_abs_sub_le_antitone}
\leanok
\uses{def:isSigmaOrderContinuousNorm}
Let $X$ be a normed vector lattice with $\sigma$-order continuous
norm, let $x \in X$, and let $(u_n)_{n \in \mathbb{N}}$,
$(v_n)_{n \in \mathbb{N}}$ be sequences in $X$ such that $v$ is
antitone with $0 \le v_n$ for all $n$ and greatest lower bound
$0$, and $\lvert u_n - x \rvert \le v_n$ for all $n$. Then $u_n$
converges to $x$ in norm.
\end{lemma}

\begin{lemma}
\label{lem:isSigmaOrderContinuousNorm-tendsto-norm-of-monotone-isLUB}
\lean{IsSigmaOrderContinuousNorm.tendsto_norm_of_monotone_isLUB}
\leanok
\uses{lem:isSigmaOrderContinuousNorm-tendsto-of-monotone-isLUB}
Let $X$ be a normed vector lattice with $\sigma$-order continuous
norm. If $(u_n)_{n \in \mathbb{N}}$ is a monotone sequence in $X$
with least upper bound $x \in X$, then $\lVert u_n \rVert$
converges to $\lVert x \rVert$ in $\mathbb{R}$.
\end{lemma}

\begin{theorem}
\label{thm:conditionallyCompleteLatticeOf-isOrderContinuousNorm}
\lean{BanachLattice.conditionallyCompleteLatticeOf_isOrderContinuousNorm}
\leanok
\uses{def:isOrderContinuousNorm, def:banach-lattice,
  thm:conditionallyCompleteLatticeOfPosSet}
Every Banach lattice with order continuous norm is Dedekind
complete: every nonempty bounded above subset admits a least
upper bound.
\end{theorem}

\section{Nakano's theorem}

Nakano's theorem characterises order continuity of the norm of a
Banach lattice in terms of monotone bounded sequences: it is
equivalent to $\sigma$-conditional completeness paired with
$\sigma$-order continuity, and equivalent to norm convergence of
every monotone bounded above sequence to its supremum.

\begin{theorem}
\label{thm:isOrderContinuousNorm-tfae}
\lean{BanachLattice.tendsto_of_monotone_bddAbove,
  BanachLattice.sigmaConditionallyCompleteLattice_of_mono_bddAbove_tendsto,
  BanachLattice.isSigmaOrderContinuousNorm_of_mono_bddAbove_tendsto,
  BanachLattice.isOrderContinuousNorm_of_isSigmaConditionallyCompleteLattice}
\leanok
\uses{def:isSigmaOrderContinuousNorm, def:isOrderContinuousNorm,
  def:banach-lattice, def:vector-lattice,
  def:sigmaConditionallyCompleteLattice,
  thm:conditionallyCompleteLatticeOf-isOrderContinuousNorm,
  thm:sigmaConditionallyCompleteLatticeOfPosSeq,
  lem:isSigmaOrderContinuousNorm-tendsto-of-monotone-isLUB,
  lem:isSigmaOrderContinuousNorm-tendsto-of-antitone-isGLB,
  lem:isOrderContinuousNorm-toIsSigmaOrderContinuousNorm,
  thm:isGLB-of-antitone-tendsto}
Let $X$ be a Banach lattice. The following are equivalent:
\begin{enumerate}
  \item the norm of $X$ is order continuous;
  \item $X$ is $\sigma$-conditionally complete and the norm of $X$
    is $\sigma$-order continuous;
  \item every monotone bounded above sequence in $X$ has a least
    upper bound towards which it converges in norm.
\end{enumerate}
\end{theorem}

\section{Meyer--Nieberg's theorem}

Meyer--Nieberg's theorem characterises order continuity of the
norm of a Banach lattice in terms of order bounded pairwise
disjoint sequences. As a consequence, an order bounded pairwise
disjoint set of nonzero elements in such a Banach lattice is at
most countable.

\begin{theorem}
\label{thm:disjoint-bddAbove-tendsto-zero}
\lean{BanachLattice.disjoint_bddAbove_tendsto_zero}
\leanok
\uses{def:isOrderContinuousNorm, def:banach-lattice,
  def:isVLDisjoint, thm:isOrderContinuousNorm-tfae}
Let $X$ be a Banach lattice with order continuous norm and let
$(u_n)_{n \in \mathbb{N}}$ be a sequence in $X$ such that $u_i$
and $u_j$ are disjoint whenever $i \ne j$ and $\{ \lvert u_n
\rvert : n \in \mathbb{N} \}$ is bounded above. Then $u_n$
converges to $0$ in norm.
\end{theorem}

\begin{theorem}
\label{thm:countable-of-pairwise-disjoint-bddAbove}
\lean{BanachLattice.countable_of_pairwise_disjoint_bddAbove}
\leanok
\uses{def:isOrderContinuousNorm, def:banach-lattice,
  def:isVLDisjoint, thm:disjoint-bddAbove-tendsto-zero}
Let $X$ be a Banach lattice with order continuous norm and let $S
\subseteq X$ be a set of pairwise disjoint nonzero elements such
that $\{ \lvert x \rvert : x \in S \}$ is bounded above. Then $S$
is at most countable.
\end{theorem}

\section{Ando's theorem}

Ando's theorem characterises order continuity of the norm of a
Banach lattice in terms of norm-closed order ideals: the norm is
order continuous if and only if every norm-closed order ideal is
a band, in which case every such ideal is in fact a projection
band.

\begin{theorem}
\label{thm:band-of-isClosed-orderIdeal}
\lean{BanachLattice.band_of_isClosed_orderIdeal}
\leanok
\uses{def:isOrderContinuousNorm, def:banach-lattice,
  def:order-ideal, def:band, lem:band-ofPosDirectedSSupMem}
Let $X$ be a Banach lattice with order continuous norm. Every
norm-closed order ideal of $X$ is a band.
\end{theorem}

\begin{theorem}
\label{thm:projectionBand-of-isClosed-orderIdeal}
\lean{BanachLattice.projectionBand_of_isClosed_orderIdeal}
\leanok
\uses{def:isOrderContinuousNorm, def:banach-lattice,
  def:order-ideal, def:projectionBand,
  thm:band-of-isClosed-orderIdeal,
  lem:closedOrderIdeal-sup, lem:isClosed-disjointComplement,
  lem:band-disjointComplement,
  thm:band-eq-disjointComplement-disjointComplement,
  lem:disjointComplement-anti,
  lem:disjointComplement-inter-eq-zero,
  thm:projectionBand-iff-add-disjointComplement,
  lem:orderIdeal-sum}
Let $X$ be a Banach lattice with order continuous norm. Every
norm-closed order ideal of $X$ is a projection band.
\end{theorem}

\section{The decomposition lemma}

Order continuity of the norm forces the principal projection
property: the principal band generated by any element is a
projection band, so the principal band projection $P_a$ along
$a \in X$ is well defined. The decomposition lemma then says
that the principal band projections along a maximal disjoint
family of strictly positive elements decompose every vector
unconditionally.

\begin{lemma}
\label{lem:hasPrincipalProjectionProperty-of-isOrderContinuousNorm}
\lean{BanachLattice.hasPrincipalProjectionProperty_of_isOrderContinuousNorm}
\leanok
\uses{def:isOrderContinuousNorm, def:banach-lattice,
  def:hasPrincipalProjectionProperty,
  thm:conditionallyCompleteLatticeOf-isOrderContinuousNorm,
  thm:hasPrincipalProjectionProperty-of-isSigmaOrderComplete}
Every Banach lattice with order continuous norm has the
principal projection property.
\end{lemma}

\begin{theorem}
\label{thm:hasSum-principalBandProjection}
\lean{BanachLattice.hasSum_principalBandProjection}
\leanok
\uses{def:isOrderContinuousNorm, def:banach-lattice,
  def:isMaximalDisjoint,
  def:band-principalBandProjection,
  lem:hasPrincipalProjectionProperty-of-isOrderContinuousNorm,
  thm:isLUB-principalBandProjection-of-isMaximalDisjoint}
Let $X$ be a Banach lattice with order continuous norm, let
$\Lambda \subseteq X$ be a maximal disjoint set whose elements
are strictly positive, and let $x \in X$. Then the family of
principal band projections $\bigl( P_a x \bigr)_{a \in \Lambda}$
is unconditionally summable with sum $x$.
\end{theorem}
